\documentclass[reqno]{amsart}
\usepackage[T1]{fontenc}
\usepackage{amsmath,amssymb,amsthm}
\usepackage{xr-hyper}
\usepackage[hidelinks]{hyperref}
\externaldocument[tt-]{20261125T004100--tate-twists-and-sha__workbook}
\externaldocument[ec-]{20261125T014500--eisenstein-congruences-and-selmer-classes__workbook}
\externaldocument[pc-]{20261125T014700--pseudocharacters-and-the-length-bound__workbook}
\externaldocument[lt-]{20261125T020400--local-terms-congruence-defects-and-selmer-windows__workbook}
\externaldocument[tw-]{20261126T020420--exceptional-twists-in-cyclotomic-towers__workbook}
\externaldocument[mi-]{20261126T020430--the-minus-part-by-gauss-sums__workbook}

\newcounter{chap}
% Chapter numbers, generated from 20261125T012000--grh-workbook__guide.org
% by scripts/make-chapters.py.  Do not edit: change the order of the
% index and rerun the script (./preview-site.sh and CI do this).
%
% Usage in a sheet, after \newcounter{chap}:
%   % Chapter numbers, generated from 20261125T012000--grh-workbook__guide.org
% by scripts/make-chapters.py.  Do not edit: change the order of the
% index and rerun the script (./preview-site.sh and CI do this).
%
% Usage in a sheet, after \newcounter{chap}:
%   % Chapter numbers, generated from 20261125T012000--grh-workbook__guide.org
% by scripts/make-chapters.py.  Do not edit: change the order of the
% index and rerun the script (./preview-site.sh and CI do this).
%
% Usage in a sheet, after \newcounter{chap}:
%   \input{chapters}\chapnumber{<sheet ID>}
\makeatletter
\providecommand\chapters@set[2]{\@namedef{chapters@#1}{#2}}
\chapters@set{20261125T001000}{1}% zeros-to-ramp-sums
\chapters@set{20261125T001100}{2}% inert-primes-and-estimate-b
\chapters@set{20261125T005000}{3}% galois-fourier-and-averaging
\chapters@set{20261125T005500}{4}% columns-from-a-sphere
\chapters@set{20261125T000000}{5}% slope-method
\chapters@set{20261126T030100}{6}% a-twist-for-each-prime
\chapters@set{20261125T003000}{7}% the-chain-and-the-square-root
\chapters@set{20261125T010000}{8}% rank-lemma-and-zero-estimate
\chapters@set{20261125T010500}{9}% poincare-series-and-near-orthogonality
\chapters@set{20261124T223000}{10}% elliptic-curves-for-the-proof
\chapters@set{20261125T009500}{11}% fock-model-and-theta-intertwiners
\chapters@set{20261124T223100}{12}% complex-multiplication-and-inert-reduction
\chapters@set{20261125T008000}{13}% one-cm-determinant
\chapters@set{20261124T214500}{14}% supersingular-theta-divisibility
\chapters@set{20261124T214600}{15}% the-isotropic-plane
\chapters@set{20261125T013700}{16}% pairings-on-a-formal-group-of-height-two
\chapters@set{20261125T013500}{17}% hecke-nondegeneracy-on-the-tree
\chapters@set{20261126T030200}{18}% the-bound-is-attained
\chapters@set{20261125T009000}{19}% kronecker-limit-and-chowla-selberg
\chapters@set{20261125T002000}{20}% definite-theta-lifts
\chapters@set{20261125T015500}{21}% siegel-weil-in-the-smallest-case
\chapters@set{20261125T002100}{22}% definite-theta-lifts-and-norms
\chapters@set{20261125T016500}{23}% hecke-operators-and-theta-lifts
\chapters@set{20261125T016700}{24}% divisibility-along-the-chain
\chapters@set{20261125T004000}{25}% denominators-and-sha
\chapters@set{20261125T004100}{26}% tate-twists-and-sha
\chapters@set{20261126T020430}{27}% the-minus-part-by-gauss-sums
\chapters@set{20261125T014100}{28}% congruence-modules
\chapters@set{20261125T014500}{29}% eisenstein-congruences-and-selmer-classes
\chapters@set{20261125T014700}{30}% pseudocharacters-and-the-length-bound
\chapters@set{20261125T011000}{31}% galois-representations-black-box
\chapters@set{20261126T030300}{32}% the-data
\chapters@set{20261126T020210}{33}% bessel-poincare-series-and-their-gram-operators
\chapters@set{20261126T030400}{34}% the-seesaw-identity
\chapters@set{20261126T030500}{35}% divisibility-of-the-composite-lift
\chapters@set{20261126T030600}{36}% theorem-a-short-forms
\chapters@set{20261126T030700}{37}% the-determinant-of-a-chain
\chapters@set{20261126T030800}{38}% the-cost-of-one-lift
\chapters@set{20261126T030900}{39}% denominators-of-projections-to-lifts
\chapters@set{20261126T031000}{40}% twists-with-few-congruences
\chapters@set{20261126T031100}{41}% theorem-c-and-the-main-inequality
\providecommand\chapnumber[1]{%
  \@ifundefined{chapters@#1}%
    {\PackageError{chapters}{Sheet #1 is not listed in the index}%
      {List it in 20261125T012000--grh-workbook__guide.org\MessageBreak
       and run scripts/make-chapters.py.}}%
    {\edef\thechap{\@nameuse{chapters@#1}}}}
\makeatother
\chapnumber{<sheet ID>}
\makeatletter
\providecommand\chapters@set[2]{\@namedef{chapters@#1}{#2}}
\chapters@set{20261125T001000}{1}% zeros-to-ramp-sums
\chapters@set{20261125T001100}{2}% inert-primes-and-estimate-b
\chapters@set{20261125T005000}{3}% galois-fourier-and-averaging
\chapters@set{20261125T005500}{4}% columns-from-a-sphere
\chapters@set{20261125T000000}{5}% slope-method
\chapters@set{20261126T030100}{6}% a-twist-for-each-prime
\chapters@set{20261125T003000}{7}% the-chain-and-the-square-root
\chapters@set{20261125T010000}{8}% rank-lemma-and-zero-estimate
\chapters@set{20261125T010500}{9}% poincare-series-and-near-orthogonality
\chapters@set{20261124T223000}{10}% elliptic-curves-for-the-proof
\chapters@set{20261125T009500}{11}% fock-model-and-theta-intertwiners
\chapters@set{20261124T223100}{12}% complex-multiplication-and-inert-reduction
\chapters@set{20261125T008000}{13}% one-cm-determinant
\chapters@set{20261124T214500}{14}% supersingular-theta-divisibility
\chapters@set{20261124T214600}{15}% the-isotropic-plane
\chapters@set{20261125T013700}{16}% pairings-on-a-formal-group-of-height-two
\chapters@set{20261125T013500}{17}% hecke-nondegeneracy-on-the-tree
\chapters@set{20261126T030200}{18}% the-bound-is-attained
\chapters@set{20261125T009000}{19}% kronecker-limit-and-chowla-selberg
\chapters@set{20261125T002000}{20}% definite-theta-lifts
\chapters@set{20261125T015500}{21}% siegel-weil-in-the-smallest-case
\chapters@set{20261125T002100}{22}% definite-theta-lifts-and-norms
\chapters@set{20261125T016500}{23}% hecke-operators-and-theta-lifts
\chapters@set{20261125T016700}{24}% divisibility-along-the-chain
\chapters@set{20261125T004000}{25}% denominators-and-sha
\chapters@set{20261125T004100}{26}% tate-twists-and-sha
\chapters@set{20261126T020430}{27}% the-minus-part-by-gauss-sums
\chapters@set{20261125T014100}{28}% congruence-modules
\chapters@set{20261125T014500}{29}% eisenstein-congruences-and-selmer-classes
\chapters@set{20261125T014700}{30}% pseudocharacters-and-the-length-bound
\chapters@set{20261125T011000}{31}% galois-representations-black-box
\chapters@set{20261126T030300}{32}% the-data
\chapters@set{20261126T020210}{33}% bessel-poincare-series-and-their-gram-operators
\chapters@set{20261126T030400}{34}% the-seesaw-identity
\chapters@set{20261126T030500}{35}% divisibility-of-the-composite-lift
\chapters@set{20261126T030600}{36}% theorem-a-short-forms
\chapters@set{20261126T030700}{37}% the-determinant-of-a-chain
\chapters@set{20261126T030800}{38}% the-cost-of-one-lift
\chapters@set{20261126T030900}{39}% denominators-of-projections-to-lifts
\chapters@set{20261126T031000}{40}% twists-with-few-congruences
\chapters@set{20261126T031100}{41}% theorem-c-and-the-main-inequality
\providecommand\chapnumber[1]{%
  \@ifundefined{chapters@#1}%
    {\PackageError{chapters}{Sheet #1 is not listed in the index}%
      {List it in 20261125T012000--grh-workbook__guide.org\MessageBreak
       and run scripts/make-chapters.py.}}%
    {\edef\thechap{\@nameuse{chapters@#1}}}}
\makeatother
\chapnumber{<sheet ID>}
\makeatletter
\providecommand\chapters@set[2]{\@namedef{chapters@#1}{#2}}
\chapters@set{20261125T001000}{1}% zeros-to-ramp-sums
\chapters@set{20261125T001100}{2}% inert-primes-and-estimate-b
\chapters@set{20261125T005000}{3}% galois-fourier-and-averaging
\chapters@set{20261125T005500}{4}% columns-from-a-sphere
\chapters@set{20261125T000000}{5}% slope-method
\chapters@set{20261126T030100}{6}% a-twist-for-each-prime
\chapters@set{20261125T003000}{7}% the-chain-and-the-square-root
\chapters@set{20261125T010000}{8}% rank-lemma-and-zero-estimate
\chapters@set{20261125T010500}{9}% poincare-series-and-near-orthogonality
\chapters@set{20261124T223000}{10}% elliptic-curves-for-the-proof
\chapters@set{20261125T009500}{11}% fock-model-and-theta-intertwiners
\chapters@set{20261124T223100}{12}% complex-multiplication-and-inert-reduction
\chapters@set{20261125T008000}{13}% one-cm-determinant
\chapters@set{20261124T214500}{14}% supersingular-theta-divisibility
\chapters@set{20261124T214600}{15}% the-isotropic-plane
\chapters@set{20261125T013700}{16}% pairings-on-a-formal-group-of-height-two
\chapters@set{20261125T013500}{17}% hecke-nondegeneracy-on-the-tree
\chapters@set{20261126T030200}{18}% the-bound-is-attained
\chapters@set{20261125T009000}{19}% kronecker-limit-and-chowla-selberg
\chapters@set{20261125T002000}{20}% definite-theta-lifts
\chapters@set{20261125T015500}{21}% siegel-weil-in-the-smallest-case
\chapters@set{20261125T002100}{22}% definite-theta-lifts-and-norms
\chapters@set{20261125T016500}{23}% hecke-operators-and-theta-lifts
\chapters@set{20261125T016700}{24}% divisibility-along-the-chain
\chapters@set{20261125T004000}{25}% denominators-and-sha
\chapters@set{20261125T004100}{26}% tate-twists-and-sha
\chapters@set{20261126T020430}{27}% the-minus-part-by-gauss-sums
\chapters@set{20261125T014100}{28}% congruence-modules
\chapters@set{20261125T014500}{29}% eisenstein-congruences-and-selmer-classes
\chapters@set{20261125T014700}{30}% pseudocharacters-and-the-length-bound
\chapters@set{20261125T011000}{31}% galois-representations-black-box
\chapters@set{20261126T030300}{32}% the-data
\chapters@set{20261126T020210}{33}% bessel-poincare-series-and-their-gram-operators
\chapters@set{20261126T030400}{34}% the-seesaw-identity
\chapters@set{20261126T030500}{35}% divisibility-of-the-composite-lift
\chapters@set{20261126T030600}{36}% theorem-a-short-forms
\chapters@set{20261126T030700}{37}% the-determinant-of-a-chain
\chapters@set{20261126T030800}{38}% the-cost-of-one-lift
\chapters@set{20261126T030900}{39}% denominators-of-projections-to-lifts
\chapters@set{20261126T031000}{40}% twists-with-few-congruences
\chapters@set{20261126T031100}{41}% theorem-c-and-the-main-inequality
\providecommand\chapnumber[1]{%
  \@ifundefined{chapters@#1}%
    {\PackageError{chapters}{Sheet #1 is not listed in the index}%
      {List it in 20261125T012000--grh-workbook__guide.org\MessageBreak
       and run scripts/make-chapters.py.}}%
    {\edef\thechap{\@nameuse{chapters@#1}}}}
\makeatother

\ifcsname chapters@20261126T020440\endcsname
  \chapnumber{20261126T020440}
\else
  \renewcommand{\thechap}{?}
\fi
\renewcommand{\thesection}{\thechap.\arabic{section}}
\theoremstyle{definition}
\newtheorem{problem}{Problem}[section]
\theoremstyle{plain}
\newtheorem{theorem}{Theorem}[section]
\newtheorem{lemma}[theorem]{Lemma}
\theoremstyle{remark}
\newtheorem*{leanremark}{In the Lean formalization}

\newcommand{\Z}{\mathbb{Z}}
\newcommand{\Q}{\mathbb{Q}}
\newcommand{\C}{\mathbb{C}}
\newcommand{\F}{\mathbb{F}}
\newcommand{\Zl}{\Z_\ell}
\newcommand{\Ql}{\Q_\ell}
\newcommand{\Qlb}{\overline{\Q}_\ell}
\newcommand{\Oc}{\mathcal{O}}
\newcommand{\Bc}{\mathcal{B}}
\newcommand{\Gal}{\operatorname{Gal}}
\newcommand{\Hom}{\operatorname{Hom}}
\newcommand{\tr}{\operatorname{tr}}
\newcommand{\len}{\operatorname{len}}
\newcommand{\Sha}{\mathrm{Sha}}
\newcommand{\tl}[1]{\texttt{#1}}
\setlength{\emergencystretch}{3em}

\title{The plus part, by congruences}
\author{}
\date{November 26, 2026}

\begin{document}
\maketitle
\refstepcounter{chap}\label{chap}

\paragraph{Prerequisites.}
Chapter~\ref{mi-chap}, Section~\ref{mi-sec:signs}: the fields \(E=E(\psi,a,q,d)\), the plus part \(\Sha^+_\ell(E,e)\) and its length \(h^+_{\ell,e}(E)\).  Chapter~\ref{tt-chap}: \(W_e\), \(\Sha_\ell(\Q,e;\eta)\) and Problem~\ref{tt-pr:caseZ}.  Chapter~\ref{ec-chap}, Section~\ref{ec-sec:shapiro}: the coefficient ring \(\Oc\), with ramification index \(e_\Oc\) over \(\Zl\).  Chapter~\ref{pc-chap}: the length bound of Theorem~\ref{pc-thm:model} and its factor two (Section~\ref{pc-sec:two}).  Chapter~\ref{tw-chap}: layers.  Galois representations attached to modular forms, Bloch--Kato Selmer groups \cite{BK}, and Bertrand's postulate \cite[Ch.~2]{AZ}.  Throughout, \(\ell\) is an odd prime, \(N=\operatorname{lcm}(2^{a+2},N_\psi)qd\) is the modulus of \(E\), and \(v_\ell\) is normalized by \(v_\ell(\ell)=1\).

\section{Characters and cases}\label{sec:cases}

On the plus part no \(L\)-value is available (Problem~\ref{mi-pr:parity}).  How can one bound it at all, uniformly in \(\ell\) and \(e\)?  The answer in the formalization is a case analysis over \((\ell,e,\eta)\), with a different arithmetic argument in each case, and in the main case an integer built from automorphic forms in place of a Bernoulli number.  First the reduction to characters.  For a character \(\eta\) modulo \(N\) with values in \(\Oc\), let \(\Sha^{(N)}_\ell(\Q,e;\eta)\) be the group of classes in \(\Sha_\ell(\Q,e;\eta)_\Oc\) that vanish also on the decomposition groups at the primes dividing \(N\).  The character \(\eta\) is \emph{plus for \(e\)} if \(\eta(-1)=(-1)^e\).

\begin{lemma}[reduction to characters]\label{lem:red}
If \(\ell\nmid[E:\Q]\) and \(e\ge1\), then \(e_\Oc\,h^+_{\ell,e}(E)\le\sum_\eta\len_\Oc\Sha^{(N)}_\ell(\Q,e;\eta)\), the sum over the characters \(\eta\) of \(\Gal(E/\Q)\) that are plus for \(e\).
\end{lemma}

\noindent\emph{Status}: proved in the Lean formalization; we have read the statement, not the proof.  Without the extra local conditions it would follow from Problems~\ref{ec-pr:sha-decomp} and~\ref{mi-pr:parts}.

\begin{problem}\label{pr:caseZ}
\(\Sha^{(N)}_\ell(\Q,e;\eta)\subseteq\Sha_\ell(\Q,e;\eta)_\Oc\).  Hence, if \(\ell\ge5\), \(e\ge2\) is even and \((\ell-1)\mid2e\), the summand of the trivial character vanishes.
\end{problem}

Let \(o\) be the order of \(\psi\).  Let \(p_K\) be the least prime \(p\equiv3\pmod4\) with \(p\ge7\) and \(p\nmid N_\psi\), let \(h_*\) be the class number of \(\Q(\sqrt{-p_K})\), and let \(p_1\) be the least prime \(p\equiv1\) modulo the odd part of \(o\) with \(p\nmid2N_\psi p_K\).  The set of \emph{bad primes} of \(\psi\) is
\[
 \Bc_\psi=\{3\}\cup\{\text{odd primes dividing }o\,N_\psi\,p_K\,h_*\,p_1\}.
\]
It depends only on \(\psi\).  The role of \(p_K\) and \(h_*\) is presumably to supply the imaginary quadratic field of Section~\ref{sec:integer}; we have not checked this, nor traced the role of \(p_1\).

\begin{problem}\label{pr:bad-primes}
For \(\psi=1\), \(\Bc_\psi=\{3,7\}\).  For the quadratic character of conductor \(5\), \(\Bc_\psi=\{3,5,7\}\).  For a character of order \(3\) and conductor \(7\), \(p_K=11\), \(h_*=1\), \(p_1=13\) and \(\Bc_\psi=\{3,7,11,13\}\).
\end{problem}

\begin{problem}\label{pr:gd-auto}
Let \(\ell\notin\Bc_\psi\) and let \(\eta\) be a character of \(\Gal(E/\Q)\).  Then \(\ell\nmid[E:\Q]\), \(\ell\nmid o(\eta)\), and the component of \(\eta\) at \(\ell\) (its restriction to \((\Z/\ell^{v_\ell(N)})^\times\)) has order at most \(2\).
\end{problem}

Let \(\chi_P(\eta)\) be the character obtained from \(\eta\) by removing its part of \(\ell\)-power order and then its component at \(\ell\).  For \(0<e<t\) and \(\eta\) plus for \(e\), the formalization distinguishes:
\begin{itemize}
\item[I] \(\ell\notin\Bc_\psi\) and \((\ell-1)\nmid2e\);
\item[Z] \(\ell\notin\Bc_\psi\), \((\ell-1)\mid2e\) and \(\eta=1\);
\item[II] \(e\ge2\), the order of \(\chi_P(\eta)\) is even, and \(\ell\in\Bc_\psi\) or \((\ell-1)\mid2e\);
\item[III] \(\ell\notin\Bc_\psi\), \((\ell-1)\mid2e\), \(\eta\ne1\) and the order of \(\eta\) is odd;
\item[IV] \(\ell\in\Bc_\psi\), \(e\ge2\) and the order of \(\chi_P(\eta)\) is odd;
\item[V] \(\ell\in\Bc_\psi\) and \(e=1\).
\end{itemize}

\begin{problem}\label{pr:exhaustive}
Suppose \(2t+1<q\) and \(2t+1<d\).  Then every triple \((\ell,e,\eta)\) as above lies in at least one of the cases I--V and Z.  In Cases II and III with \(\ell\notin\Bc_\psi\), \(\ell\le2e+1<2t+1\).
\end{problem}

Case~Z vanishes by Problem~\ref{pr:caseZ}.  Cases II and III have \(\ell\) small, Case~I has \(\ell\) arbitrary, and Cases II, IV and V with \(\ell\in\Bc_\psi\) involve finitely many \(\ell\).

\section{The congruence integer}\label{sec:integer}

In Case~I the Bernoulli number of the minus part is replaced by an integer.  Let \(K\) be an imaginary quadratic field, \(w\) a prime of \(K\) above a prime \(p\ge7\) with \(p\nmid N\Delta_K\), \(Q=\lvert\Oc_K/w\rvert\), \(\ell\ne p\), \(\sigma\in G_K\) an arithmetic Frobenius at \(w\), and \(\Psi=\eta\chi_\ell^e\), a character of \(G_\Q\) that we also restrict to \(G_K\).  Since \(\Psi\oplus1\oplus\Psi^{-1}\) is self-dual, \(\sigma\) and \(\sigma^{-1}\) have the same trace on it.  For a root of unity \(z\) put
\[
 t_z=Q^{-e}z+1+Q^ez^{-1}.
\]

\begin{problem}\label{pr:eis-trace}
The trace of \(\sigma^{-1}\) on \(\Psi\oplus1\oplus\Psi^{-1}\) is \(t_z\) with \(z=\eta(\sigma)^{-1}\), and \(Q^et_z\) is an algebraic integer.
\end{problem}

The representations that may be congruent to this Eisenstein point are three-dimensional representations \(\rho\) of \(G_K\) of two kinds, called \emph{admissible} below: those attached to automorphic representations of a unitary group in three variables with weights \((e,0,-e)\) and bounded level, and those of the form \(1\oplus W\) with \(W\) attached to a newform of weight \(2e+1\).  We have not transcribed the precise conditions.

\begin{lemma}[one character]\label{lem:one}
Let \(K\), \(\ell\), \(e\), \(N\) and \(\eta\) satisfy: \(e\ge1\), \(N\) even, \(\ell\nmid o(\eta)\), the component of \(\eta\) at \(\ell\) has order at most \(2\), \(\eta\) is plus for \(e\), \(\ell\nmid\Delta_Kh_K\), \((\Delta_K,N)=1\), \((\ell-1)\nmid2e\), and a further condition relating \(K\) and \(\eta\).  Let \(\sigma\in G_K\) and \(t=\Psi(\sigma)+1+\Psi(\sigma)^{-1}\), and suppose \(\tr\rho(\sigma)\ne t\) for every admissible \(\rho\).  Then \(\Sha^{(N)}_\ell(\Q,e;\eta)\) is finite, and there is a finite set \(X\) of pairwise non-isomorphic admissible \(\rho\) with
\[
 \len_\Oc\Sha^{(N)}_\ell(\Q,e;\eta)\le2e_\Oc\sum_{\rho\in X}v_\ell\bigl(\tr\rho(\sigma)-t\bigr).
\]
\end{lemma}

\begin{lemma}[the integer]\label{lem:integer}
For each \(K\) there are \(c\) and \(A\) with the following property.  Let \(e,N\ge1\), \(p\), \(w\), \(Q\) be as above, and \((z_o)\) a finite family of roots of unity, stable under \(\mathrm{Aut}(\C)\).  There is a nonzero integer \(M\) with \(\log\lvert M\rvert\le c\cdot\#(z_o)\cdot(N(e+1))^A\log p\) such that, for every odd \(\ell\ne p\) and all finite sets \(X_o\) of pairwise non-isomorphic admissible \(\rho\), one has \(\tr\rho(\sigma^{-1})\ne t_{z_o}\) and \(v_\ell(M)\ge\sum_o\sum_{\rho\in X_o}v_\ell\bigl(\tr\rho(\sigma^{-1})-t_{z_o}\bigr)\).
\end{lemma}

\noindent\emph{Status}: both lemmas are proved in the Lean formalization; we have read their statements, not their proofs.  The first is a bound for a Selmer group by trace differences at congruent points, of the kind modelled in Theorem~\ref{pc-thm:model}, with the factor two of Section~\ref{pc-sec:two}; its proof uses deformation rings, automorphy lifting and the input of Section~\ref{sec:kato}.  The second counts automorphic forms on a unitary group in three variables.  Zhang describes both \cite[\S5.4]{Zh}.

Here is why such an integer exists, in a model.

\begin{problem}\label{pr:purity}
If the eigenvalues of \(\rho(\sigma^{-1})\) have absolute value \(1\) under an isomorphism \(\Qlb\simeq\C\), then \(\tr\rho(\sigma^{-1})\ne t_z\).
\end{problem}

\begin{problem}\label{pr:M-model}
Let \(a_1,\dots,a_D\) be algebraic numbers, permuted by \(\mathrm{Aut}(\C)\), with \(Q^ea_s\) integral and \(\lvert a_s\rvert\le3\) in every complex embedding, and let \((z_o)\) be as in Lemma~\ref{lem:integer}.  Then \(M=\prod_o\prod_sQ^e(a_s-t_{z_o})\) is a nonzero integer, \(\log\lvert M\rvert\le D\cdot\#(z_o)\cdot(2e\log Q+1)\), and \(v_\ell(M)=\sum_{o,s}v_\ell(a_s-t_{z_o})\) for \(\ell\ne p\).
\end{problem}

\begin{problem}\label{pr:orbit}
Fix an isomorphism \(\Qlb\simeq\C\), and assume the hypotheses of Lemma~\ref{lem:one}, other than the trace condition, for every \(\eta\) in an orbit of characters under \(\mathrm{Aut}(\C)\).  Then Lemmas~\ref{lem:one} and~\ref{lem:integer}, applied to the family \((z_\eta)\) indexed by the \(\eta\) in this orbit, with \(z_\eta=\eta(\sigma)^{-1}\), give \(\sum_\eta\len_\Oc\Sha^{(N)}_\ell(\Q,e;\eta)\le2e_\Oc v_\ell(M)\) for \(\ell\ne p\).
\end{problem}

\begin{problem}\label{pr:two-primes}
Let \(p\ne p'\), \(M_p\), \(M_{p'}\) nonzero integers, and suppose that \(L_\ell\le2v_\ell(M_p)\) for odd \(\ell\ne p\) and \(L_\ell\le2v_\ell(M_{p'})\) for odd \(\ell\ne p'\).  Then \(\sum_\ell L_\ell\log\ell\le2\log\lvert M_p\rvert+2\log\lvert M_{p'}\rvert\), the sum over odd \(\ell\).
\end{problem}

\begin{problem}\label{pr:least-primes}
The two least primes \(p<p'\) with \(p\ge7\) and \(p,p'\nmid30N\Delta_K\) satisfy \(\log p'\le C(1+\log(N\lvert\Delta_K\rvert))\), with \(C\) absolute.
\end{problem}

\section{Kato's vanishing}\label{sec:kato}

At the points \(1\oplus W\) the deformation argument must not see Selmer classes of \(W\) itself.  Let \(\theta\) be the quadratic character of \(K\).

\begin{problem}\label{pr:hom-part}
Let \(V\) be a two-dimensional representation with \(\det V=\theta\).  Then \(\Hom(1,V)\oplus\Hom(V,1)\cong V\oplus(V\otimes\theta)\).
\end{problem}

\begin{problem}\label{pr:residual}
If \(\tr W\equiv\theta\Psi+\Psi^{-1}\) modulo the maximal ideal, then on \(G_K\) the trace of \(1\oplus W\) is congruent to \(t=\Psi+1+\Psi^{-1}\): the point \(1\oplus W\) is congruent to the Eisenstein point.
\end{problem}

\begin{lemma}[Kato]\label{lem:kato}
Let \(W\) be a two-dimensional \(\ell\)-adic representation of \(G_\Q\) that is a Tate twist of the representation of a newform of weight \(2e+1\), with \(\det W=\theta\), satisfying further conditions we do not transcribe.  Then the Bloch--Kato Selmer groups \(H^1_f(\Q,W)\) and \(H^1_f(\Q,W\otimes\theta)\) vanish.
\end{lemma}

\noindent\emph{Status}: Kato proved the finiteness of the Selmer groups of modular forms at the critical twists other than the central one \cite[Thm.~14.2]{Ka}, which gives this vanishing at the relevant twists \cite[\S5.4]{Zh}.  In the Lean formalization Lemma~\ref{lem:kato} is an input of the proof of Lemma~\ref{lem:one}; we have not read how it is established there.

Problem~\ref{pr:hom-part} says that the two components of an adjoint class at \(1\oplus W\) are classes for \(W\) and \(W\otimes\theta\), so Lemma~\ref{lem:kato} removes them.

\section{The plus part}\label{sec:assembly}

\begin{problem}\label{pr:small-ell}
Suppose each pair \((\ell,\eta)\) in Cases II and III with \(\ell\notin\Bc_\psi\) contributes \(\len_\Oc\Sha^{(N)}_\ell(\Q,e;\eta)\log\ell\le(\ell Ne)^C\), with \(C\ge1\); here \(e_\Oc=1\), since \(\ell\nmid[E:\Q]\) by Problem~\ref{pr:gd-auto}.  Then the sum of these contributions over \(\ell\), \(\eta\) and \(1\le e<t\) is at most \((Nt)^{5C+4}\).
\end{problem}

\begin{problem}\label{pr:outside}
Suppose \(2t+1<q\) and \(2t+1<d\).  Assume Lemma~\ref{lem:red}, with \(\Oc\) depending on \(\ell\), the bound of Problem~\ref{pr:small-ell}, and for each \(e<t\) the bound
\[
 \sum_\ell\log\ell\sum_\eta e_\Oc^{-1}\len_\Oc\Sha^{(N)}_\ell(\Q,e;\eta)\le(Nt)^C,
\]
the sum over the \(\ell\) in Case~I.  Then
\[
 \sum_{\ell\notin\Bc_\psi}\log\ell\sum_{1\le e<t}h^+_{\ell,e}(E)\le2(Nt)^{5C+4}.
\]
\end{problem}

In the formalization the bound for Case~I comes from Problems~\ref{pr:orbit}--\ref{pr:least-primes}; the bound \((\ell Ne)^C\) in Cases II and III from a statement about automorphic forms on the group \(\mathrm{GSp}_4\) of symplectic similitudes in four variables, whose content we have not traced; and Case~V from a lifting argument for unitary groups that uses the nonvanishing of a Bernoulli number modulo \(\ell\); Zhang describes unitary groups over \(\Q(i)\) here \cite[\S5.4]{Zh}.  At \(\ell\in\Bc_\psi\), Case~IV and part of Case~V are reduced to the fixed layers of Chapter~\ref{tw-chap}.  The hypothesis used there is a bound \(B\) on layers: for every \(\ell\in\Bc_\psi\) and the layer \(L\) of degree \(\ell^{j(\ell)}\) over \(\Q(\zeta_{\ell N_\psi})\),
\[
 \sum_{1\le e<t}h_{\ell,e}(L)\le B.
\]
The result is the following.

\begin{theorem}[the plus part]\label{thm:plus}
Let \(\psi\) be primitive and \(j\) a function on primes.  There are \(C\) and primes \(r_1,r_2\) such that the following holds.  Let \(E=E(\psi,a,q,d)\) with \(q,d>\max(r_1,r_2,2t+1,N_\psi,o,p_K,h_*,p_1)\), \(t\ge2\), and suppose the bound \(B\) on layers holds.  Then
\[
 \sum_\ell\log\ell\sum_{1\le e<t}h^+_{\ell,e}(E)\le(Nt)^C+CB,
\]
the sum over odd primes \(\ell\).
\end{theorem}

\noindent\emph{Status}: proved in the Lean formalization.  For \(\ell\notin\Bc_\psi\) it is Problem~\ref{pr:outside} with the case bounds above.  We have not read the treatment of \(\ell\in\Bc_\psi\), which uses a version of Lemma~\ref{lem:red} with a loss term when \(\ell\mid[E:\Q]\).  With \(B=\exp(x^{4\varepsilon/5})\) on the windows of the next chapter, Theorem~\ref{thm:plus} is the bound for \(P_+\) in part~(I) of Theorem~\ref{lt-thm:signs}.

\begin{problem}\label{pr:Q-inside}
If \(e\) is even and \(\ell\nmid[E:\Q]\), then \(h_{\ell,e}(\Q)\le h^+_{\ell,e}(E)\).
\end{problem}

\paragraph{What remains.}
The theorem is a uniform weakening of Vandiver's conjecture for the fields \(E\).  For \(\psi=1\), \(a=2\), and \(q\), \(d\) the least admissible primes above \(\max(r_1,r_2,2t+1,7)\), which are of size polynomial in \(t\) by the prime number theorem in arithmetic progressions, Problems~\ref{pr:outside} and~\ref{pr:Q-inside} bound \(\sum_{\ell\ne3,7}\log\ell\sum_{e<t,\ e\text{ even}}h_{\ell,e}(\Q)\) by a polynomial in \(t\): the kind of bound asked for at the end of Chapter~\ref{tt-chap}.  Its proof is the part of the argument that we have been able to follow least: Lemmas~\ref{lem:red}, \ref{lem:one}, \ref{lem:integer} and~\ref{lem:kato}, and the cases with \(\mathrm{GSp}_4\), are proved in the formalization by long arguments that we have not read.

\begin{leanremark}
\raggedright In \cite{Lean} Lemma~\ref{lem:red} is \tl{HV.\allowbreak{}redPrime}; the group \(\Sha^{(N)}_\ell(\Q,e;\eta)\) has length \tl{Vlen}.  The set \(\Bc_\psi\) is \tl{Bchi}, with \tl{pK}, \tl{hStar}, \tl{pOne}; the cases are \tl{Host.\allowbreak{}CaseI}, \dots, \tl{Host.\allowbreak{}CaseV} and \tl{CaseZ}, with \(\chi_P(\eta)\) the character \tl{chiP}.  Lemma~\ref{lem:one} is \tl{OneCharU\_s}, with hypotheses \tl{GD} and admissible representations \tl{Adm3U\_s} or \tl{Adm2A}; Lemma~\ref{lem:integer} is \tl{NumberU3U\_s}, with \(t_z\) the function \tl{eis3}; Problem~\ref{pr:orbit} is \tl{propGK\_\allowbreak{}of\_\allowbreak{}numberU3U\_\allowbreak{}of\_\allowbreak{}oneCharU}, and Problem~\ref{pr:two-primes} is the passage from \tl{PropGK} to \tl{CorGK}.  Lemma~\ref{lem:kato} is \tl{KatoTwoW}.  The case bounds are \tl{CountI}, \tl{CountII}, \tl{CountIII}, \tl{CountIV}, \tl{CountV}, \tl{CountZ}, the second and third from \tl{SymplecticFourImageDense}; the bound on layers is \tl{LayerLe}, and Theorem~\ref{thm:plus} is \tl{PointHV}, from \tl{sumIn} and \tl{sumOut}.
\end{leanremark}

\section*{Hints}

\subsection*{\S\ref{sec:cases}}
\begin{description}
\item[\ref{pr:caseZ}] More local conditions give a smaller group; then Problem~\ref{tt-pr:caseZ}.
\item[\ref{pr:bad-primes}] \(h(\Q(\sqrt{-7}))=h(\Q(\sqrt{-11}))=1\).
\item[\ref{pr:gd-auto}] Problem~\ref{mi-pr:degree}; the characters of \(\Gal(E/\Q)\) are generated by \(\psi\), characters of \(2\)-power conductor, and the Legendre symbols modulo \(q\) and \(d\).
\item[\ref{pr:exhaustive}] If \((\ell-1)\mid2e\), compare \(\ell\) with \(q\) and \(d\).
\end{description}

\subsection*{\S\ref{sec:integer}}
\begin{description}
\item[\ref{pr:eis-trace}] \(\chi_\ell(\sigma)=Q\).
\item[\ref{pr:purity}] \(\lvert t_z\rvert\ge Q^e-1-Q^{-e}\) and \(Q\ge7\).
\item[\ref{pr:M-model}] \(\lvert Q^e(a_s-t_z)\rvert\le Q^{2e}+4Q^e+1\le2Q^{2e}\); the product is fixed by \(\mathrm{Aut}(\C)\); Problem~\ref{pr:purity}.
\item[\ref{pr:orbit}] Take \(X_o\) to be the set \(X\) of Lemma~\ref{lem:one} for the character of index \(o\), with \(\sigma\) replaced by \(\sigma^{-1}\).
\item[\ref{pr:two-primes}] Every \(\ell\) differs from \(p\) or from \(p'\).
\item[\ref{pr:least-primes}] Bertrand's postulate gives a prime in each interval \((2^k,2^{k+1}]\), and an integer \(m\ge1\) has at most \(\log_2m\) prime factors.
\end{description}

\subsection*{\S\ref{sec:kato}}
\begin{description}
\item[\ref{pr:hom-part}] For a two-dimensional \(V\), \(V^\vee\cong V\otimes\det V^{-1}\).
\item[\ref{pr:residual}] \(\theta\) is trivial on \(G_K\).
\end{description}

\subsection*{\S\ref{sec:assembly}}
\begin{description}
\item[\ref{pr:small-ell}] At most \(3t\) primes \(\ell\), \(t\) twists and \(N\) characters; \(3t\le t^3\).
\item[\ref{pr:outside}] Problems~\ref{pr:gd-auto}, \ref{pr:exhaustive} and~\ref{pr:caseZ}; sum the Case~I bound over \(e<t\).
\item[\ref{pr:Q-inside}] The summand of the trivial character in Problem~\ref{ec-pr:sha-decomp}, and Problem~\ref{mi-pr:parts}.
\end{description}

\begin{thebibliography}{Lean}
\bibitem[AZ]{AZ} M.~Aigner and G.~M.~Ziegler, \emph{Proofs from THE BOOK}, 6th ed., Springer, 2018.
\bibitem[BK]{BK} S.~Bloch and K.~Kato, \emph{\(L\)-functions and Tamagawa numbers of motives}, in: The Grothendieck Festschrift, Vol.~I, Progr.\ Math.\ 86, Birkh\"auser, 1990, 333--400.
\bibitem[Ka]{Ka} K.~Kato, \emph{\(p\)-adic Hodge theory and values of zeta functions of modular forms}, Ast\'erisque \textbf{295} (2004), 117--290.
\bibitem[Lean]{Lean} The Lean community, \emph{riemann-hypothesis}, a Lean formalization of Alp\"oge's argument, \url{https://github.com/leanprover-community/riemann-hypothesis}.
\bibitem[Zh]{Zh} R.~Zhang, \emph{A preliminary note on Alp\"oge's proposed proof of the generalized Riemann hypothesis}, 22 November 2026, \url{https://hexagonmath.org/pdf/2611.00809v1}.
\end{thebibliography}

\end{document}
